\section{Fields}

So far in this chapter, we have worked with real numbers: all of our
coefficients, constant terms, variables, and parameters were real
numbers. But in fact, we have not used very many properties of the
real numbers, except for the fact that we can add, subtract, multiply,
and divide them, and that these operations satisfy certain expected
properties (for example, that a non-zero number divided by itself is
equal to $1$).

This opens to door to doing linear algebra using other number systems:
for example, we can do linear algebra over the rational numbers,
complex numbers, or even over some more exotic number systems. A
system of numbers that one can do linear algebra with is called a {\em
  field}.

\begin{definition}{Field}{field}
  A \textbf{field}\index{field} is a set $F$, together with two
  operations called \textbf{addition}\index{field!addition} and
  \textbf{multiplication}\index{field!multiplication}, and together
  with two distinct elements $0$ and $1$, such that addition and
  multiplication satisfy the following properties:
  \begin{itemize}
  \item[(A1)] {Commutative law of addition:} $a+b=b+a$;
  \item[(A2)] {Associative law of addition:} $(a+b)+c = a+(b+c)$;
  \item[(A3)] {Unit law of addition:} $0+a = a$;
  \item[(A4)] {Additive inverse:} for each $a\in F$, there exists an element $(-a)\in F$ such that $a+(-a)=0$;
  \item[(M1)] {Commutative law of multiplication:} $ab=ba$;
  \item[(M2)] {Associative law of multiplication:} $(ab)c=a(bc)$;
  \item[(M3)] {Unit law of multiplication:} $1a=a$;
  \item[(M4)] {Multiplicative inverse:} for each non-zero $a\in F$, there exists an element $a^{-1}\in F$ such that $aa^{-1}=1$;
  \item[(D)] {Distributive law:} $a(b+c)=ab+ac$.
  \end{itemize}
\end{definition}

Properties (A1)--(A4) are about addition, properties (M1)--(M4) are
about multiplication, and property (D) is about both addition and
multiplication. Here are some examples and non-examples of fields:

\begin{example}{Some Fields and Some Non-Fields}{fields-examples}
  \begin{enumialphparenastyle}
    \begin{enumerate}
    \item The set $\mathbb{R}$ of real numbers is a field.
    \item The set $\mathbb{Q}$ of rational numbers is a field.
    \item The set $\mathbb{Z}$ of integers satisfies all field
      properties except for (M3). It is therefore not a field.
    \item The set $\mathbb{N}=\s{0,1,2,\ldots}$ of natural numbers
      satisfies all field properties except (A3) and (M3). It is
      therefore not a field.
    \end{enumerate}
  \end{enumialphparenastyle}
\end{example}

A field doesn't have to be infinite. The following is an example of a
field with only two elements.

\begin{example}{The Integers Modulo 2}{z-mod-2}
  Consider the set of {\em bits}\index{bit} (binary digits) $\s{0,1}$.
  We can multiply them as usual, and add them almost as usual, subject
  to the alternative rule $1+1=0$ (instead of $1+1=2$). Here is a
  summary of the rules for addition and multiplication:
  \begin{equation*}
    \begin{array}{l|ll}
      +&0&1 \\\hline
      0&0&1 \\
      1&1&0
    \end{array}
    \quad
    \begin{array}{l|ll}
      \cdot&0&1 \\\hline
      0&0&0 \\
      1&0&1
    \end{array}
  \end{equation*}
  This particular alternative arithmetic is called ``arithmetic modulo
  2''.   In computer science, the addition is also called the ``logical
  exclusive or'' operation\index{logical operation!exclusive or}, and multiplication is also called the
  ``logical and'' operation\index{logical operation!and}.  For example, we can calculate like this:
  \begin{equation*}
    \begin{array}{lll}
      1\cdot((1+0)+1) + 1 &=& 1\cdot(1+1) + 1\\
                          &=& 1\cdot 0 + 1\\
                          &=& 0 + 1\\
                          &=& 1.
    \end{array}
  \end{equation*}
  The binary digits form a field $\mathbb{Z}_2 = \s{0,1}$, also
  called \textbf{the field of integers modulo 2}\index{integers modulo 2}.
\end{example}

You can convince yourself that the 9 properties of fields are all
satisfied by the integers modulo 2. This is a bit tedious, but it can
be checked by calculations.  For example, to verify (A1), we have to
check that $0+0=0+0$, $0+1=1+0$, $1+0=0+1$, and $1+1=1+1$. Perhaps the
most interesting properties are (A4) and (M4). For (A4), we can set
$(-0)=0$ and $(-1)=1$. It may be surprising that $(-1)=1$, but you can
check for yourself that $1+(-1)=1+1=0$ when calculating modulo
$2$. For (M4), we can set $1^{-1} = 1$.

When solving systems of linear equations, we only used addition,
subtraction, multiplication, and division. Therefore, we can solve
systems of equations using the elements of any field instead of real numbers.

\begin{example}{Solving a System of Equations over $\mathbb{Z}_2$}{system-z2}
  Solve the following system in the integers modulo 2:
  \begin{equation*}
    \begin{array}{r@{~}c@{~}l}
      x + y     &=& 0 \\
      x     + z &=& 1 \\
          y + z &=& 1.
    \end{array}
  \end{equation*}
\end{example}

\begin{solution}
  As usual, we write the augmented matrix of the system of equations,
  then reduce it to {\rref} using elementary row operations. The only
  difference is that we will perform all arithmetic operations modulo
  2, rather than in the real numbers. The augmented matrix is:
  \begin{equation*}
    \leftB
    \begin{array}{rrr|r}
      1 & 1 & 0 & 0 \\
      1 & 0 & 1 & 1 \\
      0 & 1 & 1 & 1 \\
    \end{array}
    \rightB.
  \end{equation*}
  The first pivot entry is the $1$ in the upper left. We use a row
  operation to create a zero below it. Note that, because we are
  working modulo 2, adding $1$ and subtracting $1$ is the same thing.
  \begin{equation*}
    \leftB
    \begin{array}{rrr|r}
      \circled{1} & 1 & 0 & 0 \\
      1 & 0 & 1 & 1 \\
      0 & 1 & 1 & 1 \\
    \end{array}
    \rightB
    \stackrel{R_2\leftarrow R_2+R_1}{\sim}
    \leftB
    \begin{array}{rrr|r}
      \circled{1} & 1 & 0 & 0 \\
      0 & 1 & 1 & 1 \\
      0 & 1 & 1 & 1 \\
    \end{array}
    \rightB.
  \end{equation*}
  The next pivot entry is in row 2 and column 2. We create a zero
  below it by subtracting row 2 from row 3, and a zero above it by
  subtracting row 2 from row 1:
  \begin{equation*}
    \leftB
    \begin{array}{rrr|r}
      \circled{1} & 1 & 0 & 0 \\
      0 & \circled{1} & 1 & 1 \\
      0 & 1 & 1 & 1 \\
    \end{array}
    \rightB
    \stackrel{R_3\leftarrow R_3-R_2}{\sim}
    \leftB
    \begin{array}{rrr|r}
      \circled{1} & 1 & 0 & 0 \\
      0 & \circled{1} & 1 & 1 \\
      0 & 0 & 0 & 0 \\
    \end{array}
    \rightB
    \stackrel{R_1\leftarrow R_1-R_2}{\sim}
    \leftB
    \begin{array}{rrr|r}
      \circled{1} & 0 & 1 & 1 \\
      0 & \circled{1} & 1 & 1 \\
      0 & 0 & 0 & 0 \\
    \end{array}
    \rightB.
  \end{equation*}
  The resulting system is in {\rref}. We can see that the system is
  consistent, because there is no row whose left-hand side is zero and
  whose right-hand side is non-zero. We also see that there are two
  pivot columns, and therefore two pivot variables, $x$ and $y$. On
  the other hand, $z$ is a free variable, so we set it equal to a
  parameter: $z=t$. Notice that this time, the parameter $t$ is not a
  real number, but an element of $\mathbb{Z}_2$. From the equation
  $x+z=1$, we get $x=1-z=1+t$. Can you guess why I have written $1+t$
  instead of $1-t$? This is because $(-1)=1$ in the integers modulo 2.
  So $1-t = 1+(-1)t = 1+1t = 1+t$. Similarly, from the equation
  $y+z=1$, we get that $y=1+t$. Therefore, the general solution to the
  system of equations is
  \begin{equation*}
    \begin{array}{r@{~}c@{~}l}
      x &=& 1+t \\
      y &=& 1+t \\
      z &=& t,
    \end{array}
  \end{equation*}
  where $t\in\s{0,1}$ is an arbitrary parameter. Recall that this
  means that each time we plug in a particular value for $t$, we get a
  solution.

  There is one difference between solving equations in the real
  numbers and solving equations in $\mathbb{Z}_2$. In the real
  numbers, a system of equations has either no solution, a unique
  solution, or infinitely many solutions. This is because when there
  is a parameter, we automatically get infinitely many solutions. By
  contrast, in $\mathbb{Z}_2$, there are only two possible values for
  the parameter $t$, namely $t=0$ and $t=1$. For $t=0$ we get the
  solution $(x,y,z) = (1,1,0)$, and for $t=1$ we get the solution
  $(x,y,z) = (0,0,1)$.  Thus, when the general solution has one
  parameter in $\mathbb{Z}_2$, there are only two solutions, instead
  of infinitely many.
\end{solution}

\begin{example}{The Integers Modulo 5}{z-mod-5}
  Consider the set $\mathbb{Z}_5=\s{0,1,2,3,4}$, called the
  \textbf{integers modulo 5}\index{integers modulo 5}. We define their
  addition and multiplication by computing the usual addition and
  multiplication, then ``reducing'' the answer modulo 5. Here,
  ``reducing'' a number means repeatedly subtracting 5 until the
  answer is between $0$ and $4$. Imagine a clock showing 5 numbers
  instead of the usual 12:
  \begin{center}
    \includegraphics[width=3cm]{figures/clock}
  \end{center}
  If we want to calculate 3 o'clock plus 4 hours, we get 2 o'clock,
  because whenever the clock reaches 5, it resets to 0. This is how
  addition and multiplication modulo 5 are defined:
  \begin{equation*}
    \begin{array}{l|lllll}
      +&0&1&2&3&4 \\\hline
      0&0&1&2&3&4 \\
      1&1&2&3&4&0 \\
      2&2&3&4&0&1 \\
      3&3&4&0&1&2 \\
      4&4&0&1&2&3 \\
    \end{array}
    \quad
    \begin{array}{l|lllll}
      \cdot&0&1&2&3&4 \\\hline
      0&0&0&0&0&0 \\
      1&0&1&2&3&4 \\
      2&0&2&4&1&3 \\
      3&0&3&1&4&2 \\
      4&0&4&3&2&1 \\
    \end{array}
  \end{equation*}
\end{example}

We note that the integers modulo 5 form a field. Most of the
properties are tedious but easy to verify. Perhaps the most
interesting of the field properties is (M4). It says that for each
non-zero element $a$, there is another element $a^{-1}$ such that
$aa^{-1}=1$.  By looking at the multiplication table, we see that
$1\cdot 1=1$, $2\cdot 3=1$, $3\cdot 2=1$, and $4\cdot 4=1$. Therefore
we can set $1^{-1}=1$, $2^{-1}=3$, $3^{-1}=2$, and $4^{-1}=4$.

\begin{example}{Division in $\mathbb{Z}_5$}{division-z5}
  What is 2 divided by 3 in $\mathbb{Z}_5$?
\end{example}

\begin{solution}
  There are no fractions in $\mathbb{Z}_5$. The key to dividing is
  this: instead of dividing by $a$, multiply by $a^{-1}$. So we have:
  \begin{equation*}
    2/3 = 2\cdot 3^{-1} = 2\cdot 2 = 4.
  \end{equation*}
  So 2 divided by 3 equals 4. This makes sense, because 4 times 3
  equals 2, when calculating modulo 5.
\end{solution}

\begin{example}{Solving a System of Equations over $\mathbb{Z}_5$}{system-z5}
  Solve the following system of linear equations over $\mathbb{Z}_5$:
  \begin{equation*}
    \begin{array}{r@{~}c@{~}l}
      2x + z &=& 1 \\
      x+4y+z &=& 3 \\
      x+2y+3z &=& 2. \\
    \end{array}
  \end{equation*}
\end{example}

\begin{solution}
  We perform the usual Gauss-Jordan algorithm on the augmented
  matrix. The only thing to keep in mind is that, instead of dividing
  a row by $a$, we should multiply it by $a^{-1}$. And of course, we
  should reduce all intermediate results modulo $5$. For example, to
  change the first pivot entry from $2$ to $1$, we multiply by
  $2^{-1}=3$, instead of dividing by $2$.
  \begin{equation*}
    \leftB
    \begin{array}{rrr|r}
      \circled{2} & 0 & 1 & 1 \\
      1 & 4 & 1 & 3 \\
      1 & 2 & 3 & 2
    \end{array}
    \rightB
    \stackrel{R_1\leftarrow 3R_1}{\sim}
    \leftB
    \begin{array}{rrr|r}
      \circled{1} & 0 & 3 & 3 \\
      1 & 4 & 1 & 3 \\
      1 & 2 & 3 & 2
    \end{array}
    \rightB
    \stackrel{R_2\leftarrow R_2-R_1}{\stackrel{R_3\leftarrow R_3-R_1}{\sim}}
    \leftB
    \begin{array}{rrr|r}
      \circled{1} & 0 & 3 & 3 \\
      0 & \circled{4} & 3 & 0 \\
      0 & 2 & 0 & 4
    \end{array}
    \rightB
    \stackrel{R_2\leftarrow 4R_2}{\sim}
    \leftB
    \begin{array}{rrr|r}
      \circled{1} & 0 & 3 & 3 \\
      0 & \circled{1} & 2 & 0 \\
      0 & 2 & 0 & 4
    \end{array}
    \rightB
  \end{equation*}
  \begin{equation*}
    {}
    \stackrel{R_3\leftarrow R_3-2R_2}{\sim}
    \leftB
    \begin{array}{rrr|r}
      \circled{1} & 0 & 3 & 3 \\
      0 & \circled{1} & 2 & 0 \\
      0 & 0 & \circled{1} & 4
    \end{array}
    \rightB
    \stackrel{R_1\leftarrow R_1-3R_3}{\stackrel{R_2\leftarrow R_2-2R_3}{\sim}}
    \leftB
    \begin{array}{rrr|r}
      \circled{1} & 0 & 0 & 1 \\
      0 & \circled{1} & 0 & 2 \\
      0 & 0 & \circled{1} & 4
    \end{array}
    \rightB.
  \end{equation*}
  The final matrix is in {\rref}, and we see that the system has the
  unique solution $(x,y,z) = (1,2,4)$. Please double-check the
  solution with respect to the original equations.
\end{solution}

\begin{example}{The Integers Modulo 6 are not a Field}{nofield-z6}
  Are the integers modulo 6 a field? 
\end{example}

\begin{solution}
  The integers modulo 6 are the set $\mathbb{Z}_6=\s{0,1,2,3,4,5}$, with the
  addition and multiplication modulo 6:
    \[ \begin{array}{l|llllll}
    +&0&1&2&3&4&5 \\\hline
    0&0&1&2&3&4&5 \\
    1&1&2&3&4&5&0 \\
    2&2&3&4&5&0&1 \\
    3&3&4&5&0&1&2 \\
    4&4&5&0&1&2&3 \\
    5&5&0&1&2&3&4 \\
  \end{array}
  \quad
  \begin{array}{l|llllll}
    +&0&1&2&3&4&5 \\\hline
    0&0&0&0&0&0&0 \\
    1&0&1&2&3&4&5 \\
    2&0&2&4&0&2&4 \\
    3&0&3&0&3&0&3 \\
    4&0&4&2&0&4&2 \\
    5&0&5&4&3&2&1 \\
  \end{array}
  \]
  Then $\mathbb{Z}_6$ satisfies all of the field axioms except
  (M3). To see why (M3) fails, let $a=2$, and note, by looking at the
  multiplication table, that there is no $b\in\mathbb{Z}_6$ such that
  $ab=1$. Therefore, $\mathbb{Z}_6$ is not a field.
\end{solution}

We conclude this section with a fact that we will not prove.

\begin{theorem}
  Let $n$ be a positive integer. Then the set
  $\mathbb{Z}_n=\s{0,\ldots,n-1}$ of integers modulo $n$, with
  addition and multiplication modulo $n$, forms a field if and only if
  $n$ is prime. Thus, for example, $\mathbb{Z}_2$, $\mathbb{Z}_3$,
  $\mathbb{Z}_5$, $\mathbb{Z}_7$, and $\mathbb{Z}_{11}$ are fields,
  whereas $\mathbb{Z}_4$, $\mathbb{Z}_6$, and $\mathbb{Z}_9$ are not.
\end{theorem}

\void{
  

\begin{example}{Arithmetic Modulo 5}{arithmetic-mod-5}
  Consider the set $\s{0,1,2,3,4}$......
\end{example}


\section{Abstract number systems in linear algebra}

As you already know, Linear Algebra deals with subjects such as matrix
multiplication, linear combinations, solutions of systems of linear
equations, and so on. It makes heavy use of addition, subtraction,
multiplication, and division of scalars (think, for example, of the
rule for multiplying matrices). 

It turns out that most of what we do in linear algebra does not rely
on the specific laws of arithmetic. Linear algebra works equally well
over ``alternative'' arithmetics. 

\begin{example}
  Consider multiplying two matrices, using arithmetic modulo 2 instead
  of the usual arithmetic.
  \[
  \left(\begin{array}{lll}
      0&1&1\\
      1&1&0\\
      0&0&1\\
    \end{array}
  \right)\cdot
  \left(\begin{array}{lll}
      1&0&0\\
      0&1&1\\
      1&0&1\\
    \end{array}
  \right)=
  \left(\begin{array}{lll}
      1&1&0\\
      1&1&1\\
      1&0&1\\
    \end{array}
  \right)
  \]
  For example, to calculate the entry in the first row and column, we
  compute
  \[ 0\cdot 1 + 1\cdot 0 + 1\cdot 1 = 1.
  \]
\end{example}

There are important applications of linear algebra over such abstract
number systems, particularly in the area of cryptography. This is the
reason we introduce the concept of a {\em field}.

\section{The field axioms}

This means we can do linear algebra taking the real numbers, the complex
numbers, or the rational numbers as the scalars.

\begin{example}
  Consider the set $\Z_2=\s{0,1}$ from Example~\ref{exa-z2}, with the
  addition and multiplication given by the rules of arithmetic
  ``modulo 2'':
  \[ \begin{array}{l|ll}
    +&0&1 \\\hline
    0&0&1 \\
    1&1&0
  \end{array}
  \quad
  \begin{array}{l|ll}
    \cdot&0&1 \\\hline
    0&0&0 \\
    1&0&1
  \end{array}
  \]
  With these operations, $\Z_2$ is a field.
\end{example}

This means we can do linear algebra over $\Z_2$. 

\begin{nproblem}
  What is subtraction in $\Z_2$?
\end{nproblem}

\begin{nproblem}
Multiply the following matrices, taking scalars in $\Z_2$. 
\[   \left(\begin{array}{lll}
      0&1&1\\
      1&1&0\\
      0&0&1\\
    \end{array}
  \right)\cdot
  \left(\begin{array}{lll}
      0&1&0\\
      1&0&0\\
      1&1&1\\
    \end{array}
  \right)
\]
Compare your answer to what you get when doing the calculation with
rational scalars. 
\end{nproblem}

\begin{nproblem}
Find the inverse of the matrix 
\[   M=\left(\begin{array}{lll}
      0&1&1\\
      1&1&0\\
      0&0&1\\
    \end{array}
  \right)
\]
using $\Z_2$ as the set of scalars. Hint: follow the usual steps of
Gaussian elimination, but use the modulo 2 operations. Compare this to
the inverse of $M$ when interpreted over the rational numbers.
\end{nproblem}

\begin{nproblem}
  Consider the set $\s{0,1}$ with the following different addition and
  multiplication rules:
  \[ \begin{array}{l|ll}
    +&0&1 \\\hline
    0&0&1 \\
    1&1&1
  \end{array}
  \quad
  \begin{array}{l|ll}
    \cdot&0&1 \\\hline
    0&0&0 \\
    1&0&1
  \end{array}
  \]
  Note that we have set $1+1=1$. Which of the nine axioms are
  satisfied? Which of the nine axioms fail, if any? Is this a field?
\end{nproblem}

\begin{example}
  The {\em integers modulo 5} are the set $\Z_5=\s{0,1,2,3,4}$, with
  the following addition and multiplication rules:
    \[ \begin{array}{l|lllll}
    +&0&1&2&3&4 \\\hline
    0&0&1&2&3&4 \\
    1&1&2&3&4&0 \\
    2&2&3&4&0&1 \\
    3&3&4&0&1&2 \\
    4&4&0&1&2&3 \\
  \end{array}
  \quad
  \begin{array}{l|lllll}
    +&0&1&2&3&4 \\\hline
    0&0&0&0&0&0 \\
    1&0&1&2&3&4 \\
    2&0&2&4&1&3 \\
    3&0&3&1&4&2 \\
    4&0&4&3&2&1 \\
  \end{array}
  \]
  This is called ``arithmetic modulo 5'', because the numbers are
  wrapped after 4: 5 is treated the same as 0, 6 is treated the same
  as 1, 7 is treated the same as 2, and so on. With these operations,
  $\Z_5$ is a field.
\end{example}

\begin{example}
  The {\em integers modulo 6}
\end{example}

\begin{example}
  More generally, for any natural number $n\geq 2$, the {\em integers
    modulo $n$} are given by $\Z_n=\s{0,1,\ldots,n-1}$, with addition
  and multiplication ``modulo $n$''. For all $n$, $\Z_n$ satisfies the
  axioms (A1)--(A4), (M1), (M2), (M4), and (D).  However, the axiom
  (M3) is only satisfied when $n$ is a prime number. It is a fact
  that $\Z_n$ is a field if and only if $n$ is prime.
\end{example}

\begin{nproblem}
  Solve the following system of linear equations with scalars in
  $\Z_5$:
  \[ 
  \]
\end{nproblem}

\section{Elementary properties of fields}

Our goal is to make arithmetic in a field look ``as much as possible''
as arithmetic in the real numbers. For this reason, it will be useful
to state some additional algebraic laws, which are consequences of the
field axioms.

\begin{proposition}[Cancellation of addition]
  For all elements $x,y,a$ of a field, if $x+a=y+a$, then $x=y$. 
\end{proposition}

\begin{proof}
  Assume $x+a=y+a$. By axiom (A3), there exists an element $b$ such
  that $a+b=0$. But then we have:
  \[ \begin{array}{llll}
    x &=& 0+x & \mbox{by (A2)} \\
    &=& x+0 & \mbox{by (A4)} \\
    &=& x+(a+b) & \mbox{by assumption on $b$} \\
    &=& (x+a)+b & \mbox{by (A1)} \\
    &=& (y+a)+b & \mbox{by assumption} \\
    &=& y+(a+b) & \mbox{by (A1)} \\
    &=& y+0 & \mbox{by assumption on $b$} \\
    &=& 0+y & \mbox{by (A1)} \\
    &=& y & \mbox{by (A2)}
  \end{array}
  \]
  Note how all four axioms of addition have been used. \eot
\end{proof}

\begin{proposition}[Cancellation of multiplication]\label{pro-cancel-m}
  For all elements $x,y,a$ of a field, if $xa=ya$ and $a\neq 0$, then
  $x=y$.
\end{proposition}

\begin{nproblem}
  Prove Prop.~\ref{pro-cancel-m}.
\end{nproblem}

\begin{proposition}\label{pro-0a}
  For all elements $a$ of field, $0a=0$.
\end{proposition}

\begin{proof}
  Using distributivity and (A2), we have $0+0a = 0a = (0+0)a = 0a+0a$,
  therefore the claim follows by cancellation.\eot
\end{proof}

\begin{proposition}
  In any field, if $ab=0$, then $a=0$ or $b=0$.
\end{proposition}

\begin{proof}
  Suppose $a$ and $b$ are elements in a field such that $ab=0$. We
  must show that $a=0$ or $b=0$. We consider two cases:

  Case 1: $a=0$. Then the conclusion holds and we are done.

  Case 2: $a\neq 0$. In this case, by (M3), there exists an element
  $c$ such that $ac=1$. We have:
  \[ \begin{array}{llll}
    b &=& 1b &\mbox{by (FM2)} \\
    &=& (ac)b &\mbox{by definition of $c$} \\
    &=& (ca)b &\mbox{by (FM4)} \\
    &=& c(ab) &\mbox{by (FM1)} \\
    &=& c0    &\mbox{by assumption $ab=0$} \\
    &=& 0     &\mbox{by Prop.~\ref{pro-0a}}
  \end{array}
  \]
  In each of the two cases, we have proved $a=0$ or $b=0$. \eot
\end{proof}

The following four propositions show that certain elements, whose
existence is guaranteed by the field axioms, are in fact unique.

\begin{proposition}
  In a field, the element $0$ is uniquely determined by axiom (A2).
\end{proposition}

\begin{proof}
  Suppose that $z$ is another element also satisfying $z+a=a$ for all
  $a$. Then $z+0=0$, by definition of $z$, but also $0+z=z$, by
  definition of $0$. Using commutativity, it follows that
  $z=0+z=z+0=0$, so there cannot be more than one zero element.\eot
\end{proof}

\begin{proposition}\label{pro-uniq-a3}
  For any element $a$ of a field, the element $b$ in axiom (A3) is
  uniquely determined.
\end{proposition}

\begin{proof}
  Let $a$ be arbitrary, and suppose that there are two elements $b$
  and $b'$ such that both $a+b=0$ and $a+b'=0$. By commutativity,
  $b+a=0=b'+a$, and by cancellation, $b=b'$. It follows that there is
  no more than one element $b$ satisfying the condition of axiom (A3).
\end{proof}

\begin{remark}
  If $a,b$ are elements such that $a+b=0$, we usually write $b=(-a)$. 
  This notation is justified by Prop.~\ref{pro-uniq-a3}.
\end{remark}

\begin{proposition}\label{pro-uniq-fm2}
  In a field, the element $1$ is uniquely determined by axiom (FM2).
\end{proposition}

\begin{nproblem}
  Prove Prop.~\ref{pro-uniq-fm2}.
\end{nproblem}

\begin{proposition}\label{pro-uniq-fm3}
  For any element $a\neq 0$ of a field, the element $b$ in axiom (FM3)
  is uniquely determined.
\end{proposition}

\begin{nproblem}
  Prove Prop.~\ref{pro-uniq-fm3}.
\end{nproblem}

The next two propositions are also useful.

\begin{proposition}
  Distributivity also holds on the right: $(b+c)a = ba + ca$.
\end{proposition}

\begin{proof}
  This is a direct consequence of (D) and (FM4).\eot
\end{proof}

\begin{proposition}
  The following hold in any field, for all $a,b$:
  \begin{enumerate}\alphalabels
  \item $-(-a) = a$,
  \item $-(ab) = (-a)b = a(-b)$,
  \item $-a = (-1)a$.
  \end{enumerate}
\end{proposition}

\begin{proof}
  (a) By definition of $(-a)$, we have $a+(-a)=0$. Also, by definition
  of $-(-a)$ (and commutativity), we have $(-(-a))+(-a) = 0$. By
  cancellation, it follows that $a=-(-a)$.

  (b) To show that $-(ab) = (-a)b$, we need to show that $(-a)b$ is
  the negative of $ab$, in other words, that $ab+(-a)b=0$. This
  follows from the axioms as follows: 
  \[ \begin{array}{llll}
    ab+(-a)b &=& (a+(-a))b &\mbox{by distributivity}\\
    &=& 0b &\mbox{by (A3)} \\
    &=& 0 &\mbox{by Prop.~\ref{pro-0a}}
  \end{array}
  \]
  The proof of $-(ab)=a(-b)$ is similar.

  (c) By (FM2) and (b), we have $-a = -(1a) = (-1)a$. \eot
\end{proof}

\end{document}

Scalar.

In this book, almost everything we do will work over an arbitrary
field. The exceptions are in Section X, where we use sin and cos to
compute angles, and in Section Y, where we will sometimes have a
compute a square root.

}
